\chapter{Research Questions}
\label{sec:research_questions}

The preceding chapter identified a limitation of evaluating causal concept
faithfulness exclusively through interventions on individual concepts: when
concepts contribute to a model's response jointly, singleton interventions may
not fully represent their causal influence. This thesis therefore develops and
evaluates FELICS, an extension of the Causal Concept Faithfulness Metric that
accounts for plausible interactions between concepts.

The following research questions define the empirical and methodological
objectives of this work. They guide the conception and modeling of FELICS and
structure the subsequent evaluation. The concrete implementation, formal
properties, and statistical procedures used to investigate them are introduced
in the later chapters.

\begin{enumerate}[
		label=\textbf{Q\arabic*},
		ref=Q\arabic*,
		wide=0pt,
		leftmargin=0pt
	]

	\item \label{rq:correlation_diff}
	      \textbf{Impact of Interaction-Presumption Strength on Causal Effect Estimates}

	      \noindent\textit{Background.}
	      The extent to which concepts are treated jointly can differ between
	      questions. If broader presumed interactions represent cases in which
	      singleton interventions are less adequate, accounting for them should have
	      a more pronounced effect on the resulting causal-effect estimates.

	      \noindent\textit{Research question.}
	      Are stronger presumptions of interaction between concepts associated with
	      larger changes in causal effect estimates when concepts are evaluated
	      jointly?

	      \noindent\textit{Hypothesis.}
	      Stronger presumed interactions are expected to be associated with larger
	      differences between causal effects estimated from singleton and
	      interaction-aware interventions.

	\item \label{rq:convergence}
	      \textbf{Compatibility with the Baseline Metric}

	      \noindent\textit{Background.}
	      FELICS is intended to extend rather than replace the Causal Concept
	      Faithfulness Metric. When no concepts are evaluated jointly, the extension
	      should retain the behavior of the original approach.

	      \noindent\textit{Research question.}
	      Does FELICS reproduce the results of the original Causal Concept
	      Faithfulness Metric when it is restricted to interventions on individual
	      concepts?

	      \noindent\textit{Hypothesis.}
	      Under singleton-only interventions, FELICS is expected to yield equivalent
	      causal-effect estimates and faithfulness scores to the original metric.

	\item \label{rq:ranking}
	      \textbf{Preservation of Within-Question Concept Rankings}

	      \noindent\textit{Background.}
	      Joint interventions measure an effect that must subsequently be allocated
	      to the concepts involved. However, values obtained from singleton and joint
	      interventions may not be directly comparable on the probability-based
	      causal-effect scale. The proposed framework therefore requires an
	      adjustment that accounts for potential systematic scale differences while
	      retaining concept-level attributions. Such an adjustment may change the
	      magnitude assigned to individual concepts and could, in principle, alter
	      their relative ordering.

	      \noindent\textit{Research question.}
	      Does the adjustment required to make causal effects from singleton and
	      joint interventions comparable preserve the within-question ranking of
	      concept-level causal effects?

	      \noindent\textit{Hypothesis.}
	      The adjustment is expected to affect the magnitudes of causal-effect
	      estimates without substantially changing which concepts are ranked as more
	      or less causally influential within a question.

	\item \label{rq:faithfulness_increase}
	      \textbf{Effect of Interaction Modeling on Faithfulness}

	      \noindent\textit{Background.}
	      The Correlated Concepts Problem can cause the original metric to
	      underestimate the causal influence of concepts whose information overlaps
	      or interacts. This may widen the discrepancy between a concept's measured
	      causal effect and the importance attributed to it in a model's explanation.

	      \noindent\textit{Research question.}
	      Does accounting for interactions between concepts yield higher measured
	      faithfulness for the evaluated models?

	      \noindent\textit{Hypothesis.}
	      By more adequately representing jointly influential concepts, FELICS is
	      expected to increase measured faithfulness relative to an evaluation based
	      solely on singleton interventions.

	\item \label{rq:motivating_example}
	      \textbf{Resolution of the Motivating Underestimation Example}

	      \noindent\textit{Background.}
	      The motivating MedQA example discussed in the preceding chapter suggests
	      that concepts can receive low causal-effect estimates when the information
	      removed from one concept remains recoverable through another. It provides a
	      concrete case in which joint intervention is expected to address the
	      limitation of singleton-based evaluation.

	      \noindent\textit{Research question.}
	      Does FELICS resolve the specific underestimation identified as motivation
	      for this thesis by allotting higher causal effect compared to a singleton-only evaluation for the two presumably correlated concepts?

	      \noindent\textit{Hypothesis.}
	      Evaluating the relevant concepts jointly is expected to produce higher
	      causal-effect estimates for them than evaluating each concept in isolation.

\end{enumerate}

The following chapter derives the requirements and central design choices of
FELICS in response to these questions. The subsequent modeling chapter then
formalizes the framework before the evaluation methodology specifies how each
question is investigated empirically.
